% ==================================================================
%  CoreCheck.tex  --  Section 3.2.2 "Core Check"
%  Loaded from main.tex.
% ==================================================================

\subsubsection{Core Check}
\label{sec:core}

A point is a core point when its $\varepsilon_n$-neighborhood, itself
included, holds at least $\mathit{MinPts}$ points. Core Check counts the
neighbors of every point from the masks of Neighbor Check. Summing the
masks of both directions over all offsets, and adding the point itself,
gives the neighbor count of every point in a single ciphertext,
\begin{equation}
  \label{eq:count}
  \hat{\mathbf{c}} = 1 + \sum_{k=1}^{k_{\max}}
    \sum_{s \in \{k,\, n-k\}} \mathsf{adj}_s ,
\end{equation}
whose slot~$i$ holds $\hat{c}_i$, the number of points within
$\varepsilon_n$ of point~$i$.

Point~$i$ is a core point when $\hat{c}_i$ is at least
$\mathit{MinPts}$, and the server decides this with one more comparison.
The threshold is not $\mathit{MinPts}$ itself but
$\mathit{MinPts} - \tfrac{1}{2}$. The counts are integers, and so is
$\mathit{MinPts}$. A point with exactly $\mathit{MinPts}$ neighbors would
therefore give $\mathsf{Sgn}$ the input $0$, which lies in the dead zone
of the approximation (Section~\ref{sec:comparison}). The threshold
$\mathit{MinPts} - \tfrac{1}{2}$ lies halfway between the last non-core
count and the first core count, so every count is at least half a unit
away from it, and the input stays outside the dead zone once the
precision is chosen as in Section~\ref{sec:correct}.

No count can exceed the number of points, so dividing the shifted count
by $n$ keeps the input in $[-1,1]$:
\begin{equation}
  \label{eq:coreinput}
  \mathbf{y} = \tfrac{1}{n} \Bigl( \hat{\mathbf{c}}
    - \bigl( \mathit{MinPts} - \tfrac{1}{2} \bigr) \Bigr),
\end{equation}
\begin{equation}
  \label{eq:core}
  \mathsf{Core\_mask} = \tfrac{1}{2}
    \bigl( 1 + \mathsf{Sgn}(\mathbf{y}) \bigr).
\end{equation}
Slot~$i$ of $\mathsf{Core\_mask}$ is $1$ when point~$i$ is a core point
and $0$ when it is not. The sign is taken in the opposite direction to
Eq.~\eqref{eq:adj}, since here it is the larger values that pass.
$\mathsf{Core\_mask}$ is the mask vector that Label Propagation uses
next. The stage needs a single call to $\mathsf{Sgn}$, and
Algorithm~\ref{alg:core} collects its steps.

\begin{algorithm}[t]
\caption{\textsf{Core\_check}}
\label{alg:core}
\begin{algorithmic}[1]
\Require adjacency vectors $\mathsf{adj}_1, \dots, \mathsf{adj}_{k_{\max}}$;
  public $\mathit{MinPts}$, $k_{\max}$, $n$
\Ensure $\mathsf{Core\_mask}$, slot $i$ equal to $1$ iff point $i$ is a
  core point
\State $\hat{\mathbf{c}} \gets 1$
\For{$k = 1$ \textbf{to} $k_{\max}$}
  \State $\mathsf{adj}_{n-k} \gets \mathsf{Rot}_{n-k}(\mathsf{adj}_k)$
  \ForAll{$s \in \{k, n-k\}$}
    \State $\hat{\mathbf{c}} \gets \hat{\mathbf{c}} + \mathsf{adj}_s$
  \EndFor
\EndFor
\State $\mathbf{y} \gets \tfrac{1}{n} \bigl( \hat{\mathbf{c}}
  - (\mathit{MinPts} - \tfrac{1}{2}) \bigr)$
\State $\mathsf{Core\_mask} \gets
  \tfrac{1}{2}\bigl( 1 + \mathsf{Sgn}(\mathbf{y}) \bigr)$
\State \Return $\mathsf{Core\_mask}$
\end{algorithmic}
\end{algorithm}
